\label{ch:spectral}

\section{Représentation opératorielle des caractères de Dirichlet}

Les constantes de ce chapitre sont regroupées parce qu'elles partagent une construction
spectrale indépendante de leurs développements décimaux et exempte de valeurs
cibles. Soit

\[
Ne_n=ne_n,\qquad n\in\N,
\]

dans \(\ell^2(\N)\). Un caractère résiduel \(\chi\) définit l'opérateur diagonal \(Q_\chi e_n=\chi(n)e_n\) ; pourvu que la puissance soit de classe trace,

\begin{equation}
\Tr(N^{-s}Q_\chi)=L(s,\chi).
\label{eq:L-trace}
\end{equation}

L'équation conserve le module, l'orientation et la provenance de chaque
caractère ; par conséquent, elle n'identifie pas des objets résiduels distincts.

\section{Construction HMT des caractères résiduels}

L'opérateur diagonal précédent ne constitue pas la donnée algébrique initiale. La
chaîne de sélection est

\[
\begin{aligned}
\APP&\longrightarrow\TRIT\longrightarrow\TPK
\longrightarrow\{\text{résidu modulo }3,\ C_P\}\\
&\longrightarrow\{\chi_{-3},\ J=C_P\bmod3\}
\longrightarrow\{\chi_{-3},\ J^2=-I,\chi_{-4}\}\\
&\longrightarrow Q_3,Q_4.
\end{aligned}
\]

La fonctionnelle résiduelle triadique détermine la différence entre les deux classes
orientées non nulles. Dans le secteur d'incidence, on emploie la matrice de conférence
symétrique de Paley d'ordre six,
\[
C_P^{\mathsf T}=C_P,\qquad C_PC_P^{\mathsf T}=5I_6.
\]
Sa réduction \(J=C_P\bmod3\) satisfait
\[
J^2=2I_6=-I_6
\quad\text{dans }\operatorname{End}_{\mathbb F_3}(\mathbb F_3^6).
\]
Le quart de torsion détermine alors la phase cyclique
\(1,\ii,-1,-\ii\). L'application du théorème chinois des restes conserve
les deux données résiduelles et produit \(Q_{12}=Q_3Q_4\). La trace
\(\Tr(N^{-s}Q)\) est évaluée à l'étape terminale ; ni la série de Dirichlet
ni son expression décimale n'interviennent dans la sélection du caractère.

La construction est explicitée par les tableaux
\begin{equation}
\chi_{-3}(n)=
\begin{cases}
0,&3\mid n,\\
1,&n\equiv1\pmod3,\\
-1,&n\equiv2\pmod3,
\end{cases}
\qquad
\chi_{-4}(n)=
\begin{cases}
0,&2\mid n,\\
1,&n\equiv1\pmod4,\\
-1,&n\equiv3\pmod4.
\end{cases}
\label{eq:characters34-tables}
\end{equation}
En \(\ell^2(\N)\),
\[
Q_3e_n=\chi_{-3}(n)e_n,
\qquad Q_4e_n=\chi_{-4}(n)e_n.
\]
Comme ils sont diagonaux, ils commutent ; le théorème chinois des restes donne
\begin{equation}
(Q_3Q_4)e_n=\chi_{-3}(n)\chi_{-4}(n)e_n
=\chi_{12}(n)e_n.
\label{eq:q12-product-proof}
\end{equation}
Le produit conserve simultanément le résidu triadique et le quart de
torsion. Par conséquent, le caractère modulo douze est déterminé
avant d'effectuer toute évaluation de ses valeurs \(L\).

\begin{remark}[Contrôle négatif]
La chaîne est invalidée si \(Q_3\) ne distingue pas les deux classes orientées,
si \(J^2\ne-I\), si \(Q_3Q_4\) n'a pas le tableau CRT déclaré ou si
l'opérateur est sélectionné à partir de la valeur de la constante.
\end{remark}

\section{Structure complexe interne et constante de Catalan}

La structure complexe interne satisfait

\begin{equation}
J^2=-I.
\label{eq:J-square}
\end{equation}

L'égalité \(J^2=-I\) définit une représentation
\(\rho:C_4\to\operatorname{Aut}_{\mathbb F_3}(\mathbb F_3^6)\),
\(\rho(j)=J\). Pour la réaliser spectralement, on prend le caractère unitaire
\(\upsilon:C_4\to U(1)\), \(\upsilon(j)=\ii\), et l'application de phase
\(n\mapsto j^n\). Sur \(\ell^2(\N)\), la représentation diagonale induite
est
\[
U_4e_n=\upsilon(j^n)e_n=\ii^ne_n.
\]
Sa partie orientée est
\[
Q_4=\frac{U_4-U_4^*}{2\ii},\qquad
Q_4e_n=\sin\!\left(\frac{\pi n}{2}\right)e_n
=\chi_{-4}(n)e_n.
\]
C'est le relèvement explicite entre le générateur fini d'ordre quatre et l'opérateur
résiduel ; on n'identifie pas \(\mathbb F_9\) à \(\mathbb C\).

Sur les entiers impairs, le même caractère peut s'écrire

\[
\chi_J(n)=\ii^{n-1}=\chi_{-4}(n).
\]

De manière équivalente,

\begin{equation}
Q_4=\frac{U_4-U_4^*}{2\ii},
\qquad Q_4e_n=\chi_{-4}(n)e_n.
\label{eq:Q4}
\end{equation}

La constante de Catalan apparaît comme

\begin{equation}
G=\Tr(N^{-2}Q_4)=L(2,\chi_{-4})
=\sum_{n\ge0}\frac{(-1)^n}{(2n+1)^2}.
\label{eq:catalan}
\end{equation}

Le relèvement unitaire du quart de torsion engendre exactement le
caractère dont le moment quadratique est la constante de Catalan.

\begin{proof}[Preuve opératorielle]
L'opérateur \(N^{-2}Q_4\) est de classe trace parce que
\(\sum_{n\ge1}|\chi_{-4}(n)|n^{-2}<\infty\). Sa trace est donc la
somme des éléments diagonaux. Le tableau \cref{eq:characters34-tables}
annule les pairs et alterne les impairs :
\[
\Tr(N^{-2}Q_4)
=1-\frac1{3^2}+\frac1{5^2}-\frac1{7^2}+\cdots.
\]
L'identité avec \(L(2,\chi_{-4})\) est la définition de la série de
Dirichlet ; l'égalité avec \(\operatorname{Im}\operatorname{Li}_2(\ii)\)
s'obtient en prenant la partie imaginaire de
\(\sum_{n\ge1}\ii^n/n^2\).
\end{proof}

La valeur terminale est
\[
G=0.9159655941772190150546035149\ldots.
\]
La structure déterminante est l'identité \(J^2=-I\), qui induit
exactement le signe alterné des classes impaires ; le développement décimal
est son évaluation terminale.

\section{Composition résiduelle triadique--quaternaire}

La fonctionnelle résiduelle modulo trois fournit \(\chi_{-3}\). Son produit
CRT avec \(\chi_{-4}\) définit

\begin{equation}
\chi_{12}=\chi_{-3}\chi_{-4}.
\label{eq:chi12}
\end{equation}

Sur les unités modulo douze,

\[
\chi_{12}(1)=1,\quad
\chi_{12}(5)=-1,\quad
\chi_{12}(7)=-1,\quad
\chi_{12}(11)=1,
\]

et s'annule en dehors de celles-ci. Deux évaluations exactes sont

\begin{align}
L(1,\chi_{12})&=\frac{\log(2+\sqrt3)}{\sqrt3},
\label{eq:L1chi12}\\
L(2,\chi_{12})&=\frac{\pi^2}{6\sqrt3}.
\label{eq:L2chi12}
\end{align}

La première contient l'unité hyperbolique

\begin{equation}
\sqrt3L(1,\chi_{12})=\arcosh2=\log(2+\sqrt3).
\label{eq:pell2}
\end{equation}

Avec \(\log(3+2\sqrt2)=2\log(1+\sqrt2)\) et
\(\log(30+\sqrt{899})\), cette identité détermine une famille d'échelles
de Pell dont les membres conservent leurs provenances algébriques respectives.

Les évaluations terminales certifiées sont

\[
L(1,\chi_{12})=0.7603459963009463475310942549\ldots,
\]

\[
L(2,\chi_{12})=0.9497031262940093952634984917\ldots.
\]

\begin{proof}[Évaluations dodécaphasiques]
En regroupant par classes résiduelles,
\begin{align*}
L(1,\chi_{12})
&=\sum_{n\ge0}\left(
\frac1{12n+1}-\frac1{12n+5}
-\frac1{12n+7}+\frac1{12n+11}\right)\\
&=\int_0^1
\frac{1-x^4-x^6+x^{10}}{1-x^{12}}\,dx
=\int_0^1\frac{1-x^4}{1+x^6}\,dx.
\end{align*}
Une décomposition réelle dans les facteurs quadratiques de \(1+x^6\)
produit
\(\log(2+\sqrt3)/\sqrt3\). Pour le deuxième moment, si
\(\psi_1\) est la fonction trigamma,
\begin{align*}
L(2,\chi_{12})
&=\frac1{144}\bigl[
\psi_1(1/12)-\psi_1(5/12)
-\psi_1(7/12)+\psi_1(11/12)\bigr]\\
&=\frac{\pi^2}{144}\left(
\csc^2\frac\pi{12}-\csc^2\frac{5\pi}{12}\right)
=\frac{\pi^2}{6\sqrt3},
\end{align*}
où l'on a utilisé
\(\psi_1(z)+\psi_1(1-z)=\pi^2\csc^2(\pi z)\). Les deux preuves emploient
le tableau CRT, et non une valeur décimale cible.
\end{proof}

\section{Constantes d'Euler et de Stieltjes et transformée thêta--Mellin}

Soit
\[
\vartheta(t)=\sum_{n\in\Z}e^{-\pi n^2t},\qquad t>0.
\]
Pour \(\Re s>1\), la fonctionnelle thêta--Mellin retrouve l'identité convergente

\[
\pi^{-s/2}\Gamma\!\left(\frac{s}{2}\right)\zeta(s)
=\frac12\int_0^\infty
\bigl(\vartheta(t)-1\bigr)t^{s/2-1}\,dt.
\]

Le prolongement méromorphe s'obtient en scindant l'intégrale en \(t=1\), en utilisant
\(\vartheta(t)=t^{-1/2}\vartheta(1/t)\) et en séparant les termes polaires. Son développement en \(s=1\),

\begin{equation}
\zeta(s)=\frac1{s-1}+\sum_{j\ge0}
\frac{(-1)^j\gamma_j}{j!}(s-1)^j,
\label{eq:stieltjes}
\end{equation}

situe \(\gamma=\gamma_0\) et les constantes de Stieltjes dans une même
famille de moments renormalisés. Cette organisation exprime une
provenance analytique commune, sans attribuer une interprétation physique
identique à tous ses niveaux.

